% generated by roam.py handbook export -- не редактировать, править заметки
% 20260918160713-выпуклость_линии_уровня_для_инфимума_функции.org (body 3cd405f5e09a)

\begin{theorem}[{Выпуклость линии уровня для инфимума функции}]\label[theorem]{note:D053F5F9-AF9F-4517-A434-44B68A40039F}
  Если $f: E \rightarrow \R$~--- \hyperref[note:1CA14F2C-637F-41DE-9E97-72075D19E97F]{выпуклая}, $m = \inf_{x \in E} f(x)$
  Тогда $f_m$ (\hyperref[note:BF27A4AF-11BF-4A0D-85DF-5F7790BD58A9]{линия уровня})~--- \hyperref[note:316B5BFA-8B9F-4267-9086-D6593AAB5F9B]{выпуклое множество}.
\end{theorem}

\begin{remark}
  Если инфимум не достигается, то $f_m = \varnothing$~--- тоже выпуклое множество.
\end{remark}

\begin{proof}
  Мы знаем, что множество подуровня выпуклой функции~--- выпукло (см. \hyperref[note:1CD538BF-6F40-472A-A670-D38F093E4DF1]{Выпуклость множества подуровня выпуклой функции})

  А так как $\mathrm{Lev}_f(c) = \left\{ x \in E\ |\ f(x) \le c  \right\}$, и $f(x) \ge m\ \forall x$ ($m$~--- инфимум)
  Получим
  \[
  \mathrm{Lev}_f(m) = \left\{ x  \in E\ |\ f(x) \le m \right\} = \left\{ x \in E\ |\ f(x) = m \right\} = f_m
  \]
  Получили в точности \hyperref[note:BF27A4AF-11BF-4A0D-85DF-5F7790BD58A9]{линию уровня функции}
\end{proof}
